% Generated draft: requires mathematical and source review.
\documentclass[11pt]{article}
\usepackage[margin=25mm]{geometry}
\usepackage{fontspec}
\setmainfont{Latin Modern Roman}
\usepackage{amsmath,amssymb,mathtools,unicode-math}
\setmathfont{Latin Modern Math}
\usepackage{enumitem,needspace,xurl,hyperref}
\setlength{\parindent}{0pt}
\setlength{\parskip}{5pt}
\setlength{\emergencystretch}{3em}
\newcommand{\N}{\mathbb{N}}
\newcommand{\Z}{\mathbb{Z}}
\newcommand{\Q}{\mathbb{Q}}
\newcommand{\R}{\mathbb{R}}
\newcommand{\C}{\mathbb{C}}
\newcommand{\F}{\mathbb{F}}
\newcommand{\A}{\mathbb{A}}
\newcommand{\PP}{\mathbb{P}}
\newcommand{\E}{\mathbb{E}}
\newcommand{\eps}{\varepsilon}
\newcommand{\dd}{\,\mathrm{d}}
\begin{document}
\section*{UnsolvedMath ID: 30000001}
\textbf{Source problem code:} catalogue entry
\subsection{Definitions and mathematical statement}
Let \(k\) be a field, \(X/k\) a smooth projective variety of pure dimension \(d\), and \(H\) a chosen Weil cohomology theory with characteristic-zero coefficient field \(F\). For non-pure disconnected schemes, the dimension and correspondence conventions must be applied componentwise. Write
\[
 H^{\mathrm{ev}}(X)=\bigoplus_{i\text{ even}}H^i(X),\qquad
 H^{\mathrm{odd}}(X)=\bigoplus_{i\text{ odd}}H^i(X).
\]
The canonical parity projectors act as the identity on one summand and zero on the other.

In the standard homological formulation of the sign conjecture, the assertion is that these two projectors are induced by algebraic correspondences with rational coefficients. Namely, in
\[
 A_H^d(X\times X)_\Q
 =\{\text{codimension-}d\text{ rational cycles}\}/\text{homological equivalence for }H,
\]
there should exist \(e_{\mathrm{ev}},e_{\mathrm{odd}}\) with
\[
 e_{\mathrm{ev}}+e_{\mathrm{odd}}=[\Delta_X],\quad
 e_{\mathrm{ev}}^2=e_{\mathrm{ev}},\quad e_{\mathrm{odd}}^2=e_{\mathrm{odd}},\quad
 e_{\mathrm{ev}}e_{\mathrm{odd}}=e_{\mathrm{odd}}e_{\mathrm{ev}}=0,
\]
where products are composition of correspondences, and
\[
 (e_{\mathrm{ev}})_*|_{H^i(X)}=
 \begin{cases}\operatorname{id},&i\text{ even},\\0,&i\text{ odd},\end{cases}
 \qquad(e_{\mathrm{odd}})_*=\operatorname{id}-(e_{\mathrm{ev}})_*.
\]
For a correspondence \(e\), its action is \((p_2)_*(p_1^*(-)\cup\operatorname{cl}_H(e))\). The quantifier in the source is every smooth projective variety over a field; the cohomology theory and coefficient conventions must be fixed to interpret this statement.

The original instead says ``modulo an adequate equivalence relation'' without naming it. If that relation is finer than homological equivalence and supports realization, demanding actual orthogonal idempotents in its correspondence quotient includes an additional lifting issue. Algebraicity of the cohomological projectors alone does not supply Chow idempotents. For a coarser relation, the cohomological action need not descend without further hypotheses. Thus the displayed standard homological formulation records the sign-conjecture interpretation while leaving the original adequate-equivalence ambiguity explicit. Only even/odd parity is requested, not algebraic projectors for every individual cohomological degree.

The supplied literature triage dated 21 August 2026 describes the sign conjecture as open in general and reports known cases involving Shimura varieties. The exact Shimura classes, hypotheses, and projective models are not provided and are not generalized to all such varieties here. These are supplied dated assessments, not independently verified mathematical status.
\subsection{Short English statement}
The source asks for algebraic correspondences projecting to even and odd cohomology. The standard sign conjecture gives the homological version; the unspecified adequate equivalence may require extra realization or idempotent-lifting hypotheses. The supplied curation lists the general problem as open.
\subsection{Sources}
\begin{itemize}
\item Local UnsolvedMath registry, permanent ID \(30000001\), exact source code \texttt{catalogue entry}, Oberwolfach Reports Open Problems collection. Statement status \texttt{corrected\_verified}; supplied verification says only variable notation was restored.
\item Supplied report citation: ``Motives and Homotopy Theory of Schemes (2005),'' Oberwolfach Reports, pp. 785--828, DOI \texttt{10.4171/owr/2004/15}: \url{https://doi.org/10.4171/owr/2004/15}. The proposed year and year tag are 2004, whereas the citation label says 2005; bibliographic year conventions require checking.
\item Supplied literature source ``standard sign conjecture for Shimura varieties'': \url{https://arxiv.org/abs/1408.0461}. No authors or complete citation are supplied in the record.
\item Supplied ``SLMath seminar record'': \url{https://www.slmath.org/seminars/21049}. The record gives no further bibliographic details.
\item Literature checked 21 August 2026: supplied status \texttt{open}, general sign conjecture open with known Shimura-variety cases. Registry created 25 August 2026 and updated 16 September 2026. This is dated curation triage requiring independent verification before citation; no source visit is asserted.
\item Registry curation attribution: Ulam AI / UnsolvedMath Contributors, CC BY 4.0; the underlying report and literature are identified by their supplied provenance. This statement-only draft requires human mathematical and bibliographic review.
\end{itemize}
\end{document}
